%!TEX root = ../main.tex
% =============================================================
%  CAPÍTULO 7 — PRESENTACIÓN DE RESULTADOS
%                                       (Entrega 7 · 5 % · 13/10)
% =============================================================

\chapter{Presentación de resultados}\label{cap:resultados}

\begin{guia}[Guía de la cátedra: qué se evalúa en este capítulo]
\begin{enumerate}
    \item \textbf{Requerimientos formales}: \gls{rf} y \gls{rnf} con ID,
          descripción, prioridad y criterio de aceptación medible,
          trazados a los indicadores del capítulo~\ref{cap:metodologia}.
    \item \textbf{Evidencia de implementación}: capturas reales del sistema
          funcionando, enlaces al repositorio (anexo~\ref{anx:repositorios})
          y al despliegue.
    \item \textbf{Resultados cuantitativos honestos}: diseño experimental,
          métricas, matrices de confusión o su equivalente en el dominio,
          incluyendo lo que salió mal y su análisis. Reportar también lo
          malo y argumentarlo vale más que esconderlo. Hiperparámetros y
          versiones en el apéndice~\ref{apx:tecnicos}.
    \item \textbf{Validación con usuarios y expertos}: idealmente
          triangulada (métricas técnicas, encuesta, entrevista) y
          traducida a decisiones de diseño concretas. Instrumentos en los
          apéndices \ref{apx:encuesta} y \ref{apx:entrevista}.
\end{enumerate}
El sistema tiene que existir y andar: la entrega final incluye el
proyecto funcional.
\end{guia}

Este capítulo presenta la evidencia obtenida durante el desarrollo y la
validación de GraphSAST. Se organiza en cinco partes: los requerimientos
que el sistema debía cumplir, el proceso con el que se construyó, la
evidencia de que funciona, los resultados de las pruebas técnicas y los
de la validación con usuarios. Se informan tanto los resultados
favorables como los desfavorables, y la lectura de cada hipótesis se
reserva para el capítulo~\ref{cap:conclusiones}.

Los objetivos específicos del capítulo~\ref{cap:introduccion} se
identifican en adelante como OE-1 a OE-6, en el orden en que allí se
enuncian: analizador sintáctico (OE-1), construcción del grafo (OE-2),
motor de contaminación (OE-3), persistencia en Neo4j (OE-4), interfaz de
visualización (OE-5) y validación (OE-6).

% -------------------------------------------------------------
\section{Requerimientos del sistema}\label{sec:requerimientos}
% -------------------------------------------------------------

Los requerimientos se derivan de los objetivos específicos y de los
indicadores de evaluación definidos en el capítulo~\ref{cap:metodologia}
(tabla~\ref{tab:indicadores}). Cada requerimiento funcional indica el
objetivo al que responde y un criterio de aceptación verificable sobre
el sistema construido; la verificación de cada criterio se presenta en
la sección~\ref{sec:pruebas}.

\begin{table}[H]
    \centering
    \caption{Requerimientos funcionales}
    \label{tab:rf}
    \footnotesize
    \begin{tabularx}{\textwidth}{@{}l >{\hsize=.75\hsize}L >{\hsize=1.25\hsize}L l l L@{}}
        \toprule
        \tb{ID} & \tb{Nombre} & \tb{Descripción} & \tb{Prior.} & \tb{Obj.} & \tb{Criterio de aceptación} \\
        \midrule
        RF-01 & Análisis sintáctico & Obtener el árbol de sintaxis abstracta de archivos TypeScript y JavaScript & Alta & OE-1 & Los 43 casos del banco sintético se analizan sin errores \\
        RF-02 & Construcción del grafo & Generar los cinco tipos de nodo y las cinco relaciones del modelo & Alta & OE-2 & Pasan todas las pruebas de representación intermedia y flujo de datos \\
        RF-03 & Detección de caminos & Reportar caminos de un origen a un destino sensible sin saneamiento, según el catálogo JSON & Alta & OE-3 & Se detectan los 29 casos vulnerables del banco sintético \\
        RF-04 & Descarte de saneados & No reportar caminos que atraviesan un saneamiento del catálogo & Alta & OE-3 & Ninguno de los 14 casos seguros genera hallazgos \\
        RF-05 & Análisis entre archivos & Seguir el dato hacia funciones importadas desde otro archivo del proyecto & Alta & OE-3 & En los casos A y B de la validación con usuarios, el hallazgo cruza dos archivos \\
        RF-06 & Visualización & Dibujar el grafo, resaltar el camino y vincular cada nodo con su archivo y línea & Alta & OE-5 & El camino se resalta y cada nodo muestra su ubicación \\
        RF-07 & Carga de proyectos & Analizar en el navegador una carpeta e informar los archivos excluidos & Media & OE-5 & El análisis de una carpeta se completa en la versión publicada \\
        RF-08 & Reportes & Emitir los hallazgos en texto, JSON y SARIF 2.1.0 & Media & OE-6 & El SARIF valida contra el esquema oficial \\
        RF-09 & Integración continua & Devolver un código de salida distinto de cero ante hallazgos & Media & OE-6 & Pasan todas las pruebas de códigos de salida \\
        RF-10 & Persistencia en Neo4j & Guardar el grafo y resolver la detección con Cypher & Media & OE-4 & Sobre el banco sintético, Cypher y el recorrido en memoria dan los mismos hallazgos \\
        RF-11 & API local & Exponer el análisis en una API restringida a \texttt{127.0.0.1} & Baja & OE-5 & Se rechazan las rutas fuera de la carpeta habilitada \\
        \bottomrule
    \end{tabularx}
\end{table}

Los requerimientos no funcionales se expresan como umbrales sobre los
indicadores de la tabla~\ref{tab:indicadores}. El umbral de RNF-02
proviene del seguimiento del riesgo R-01 fijado en el
capítulo~\ref{cap:riesgos}, y el de RNF-07, del riesgo R-05. Los
umbrales de RNF-03 a RNF-06 se fijan en este capítulo, dado que el
capítulo~\ref{cap:metodologia} definió la forma de medición pero no un
valor objetivo.

\begin{table}[H]
    \centering
    \caption{Requerimientos no funcionales}
    \label{tab:rnf}
    \small
    \begin{tabularx}{\textwidth}{l L L l L}
        \toprule
        \tb{ID} & \tb{Atributo} & \tb{Descripción} & \tb{Prior.} & \tb{Umbral medible} \\
        \midrule
        RNF-01 & Corrección sobre casos conocidos & El motor no introduce falsos positivos ni falsos negativos sobre el banco sintético & Alta & FP = 0 y FN = 0 en los 43 casos (precisión y exhaustividad) \\
        RNF-02 & Exhaustividad sobre código real & Detección de las vulnerabilidades etiquetadas en una aplicación deliberadamente vulnerable & Alta & Exhaustividad $\geq$ 50\,\% sobre OWASP NodeGoat (riesgo R-01) \\
        RNF-03 & Rendimiento & Tiempo de análisis en la línea de comandos & Media & $\leq$ 1 s por cada 1.000 líneas (indicador de tiempo de análisis) \\
        RNF-04 & Confidencialidad & El código analizado no sale del equipo del usuario & Alta & 0 solicitudes a otros dominios durante el análisis en la versión web \\
        RNF-05 & Portabilidad & La versión web funciona sin instalación & Media & El análisis se completa en un navegador actualizado sin complementos \\
        RNF-06 & Mantenibilidad del catálogo & Las reglas se modifican sin cambiar el motor & Media & Agregar un destino o un saneamiento requiere editar solo un archivo JSON \\
        RNF-07 & Seguridad de las dependencias & El repositorio no publica dependencias con vulnerabilidades graves & Alta & 0 alertas críticas o altas abiertas al 3 de noviembre (riesgo R-05) \\
        \bottomrule
    \end{tabularx}
\end{table}

% -------------------------------------------------------------
\section{Proceso de desarrollo}\label{sec:proceso}
% -------------------------------------------------------------

El sistema se construyó en forma incremental, siguiendo el orden de
etapas previsto en el capítulo~\ref{cap:metodologia}
(tabla~\ref{tab:hitos}): primero el análisis sintáctico y la
representación intermedia, después el grafo de llamadas y el flujo de
datos, a continuación el motor de detección con su persistencia opcional
y, por último, la visualización. Cada etapa se cerró con pruebas
automatizadas antes de pasar a la siguiente. La
tabla~\ref{tab:hitos-reales} registra cuándo se completó cada etapa,
según los documentos de planificación y validación del repositorio y el
historial de cambios.

\begin{table}[H]
    \centering
    \caption{Etapas de desarrollo completadas}
    \label{tab:hitos-reales}
    \small
    \begin{tabularx}{\textwidth}{l L l}
        \toprule
        \tb{Etapa} & \tb{Resultado} & \tb{Registro} \\
        \midrule
        Hitos 1 a 3b & Análisis sintáctico, representación intermedia, grafo de llamadas y flujo de datos intra e inter-procedural & 12/06/2026 \\
        Validación temprana & Verificación de que el grafo contiene los caminos de origen a destino necesarios para la detección & 12/06/2026 \\
        Hito 5 & Motor de detección en memoria, persistencia opcional en Neo4j, catálogo CWE e informes & 02/07/2026 \\
        Banco sintético & Revisión del corpus de 43 casos y de la metodología de medición & 01/09/2026 \\
        Cuarta familia & Separación de la inyección NoSQL (CWE-943) de la inyección SQL & 15/09/2026 \\
        Línea de comandos & Exclusiones, códigos de salida, SARIF, API local y análisis entre archivos & 22/09/2026 \\
        Versión web & Análisis en el navegador con carga de carpetas y publicación en Vercel & 23/09/2026 \\
        Código real & Validación sobre vulnerabilidades documentadas del ecosistema npm & 23/09/2026 \\
        \bottomrule
    \end{tabularx}
\end{table}

El desarrollo se apartó del plan en dos aspectos. El primero es el
alcance: el análisis entre archivos, la versión que se ejecuta en el
navegador y la familia CWE-943 no figuraban en el alcance del
capítulo~\ref{cap:metodologia} y se incorporaron durante el desarrollo,
lo que dio origen al riesgo R-03 del capítulo~\ref{cap:riesgos}. El
segundo es el modo de ejecución: el diseño preveía la base de datos de
grafos como motor principal, y el sistema final resuelve la detección por
defecto con un recorrido en memoria, dejando Neo4j como alternativa
(capítulo~\ref{cap:marco-tecnologico}). Desde el hito 5 el motor sobre
Neo4j es opcional porque requiere una instancia activa, y en la versión
web, que se ejecuta en el navegador, el recorrido en memoria es el único
disponible.

El diseño del sistema, su organización en paquetes y el modelo de datos
del grafo se describen en el capítulo~\ref{cap:marco-tecnologico}. La
integración entre el motor y la interfaz se resolvió con un punto de
entrada del motor sin dependencias de Node.js, que la versión web ejecuta
en un proceso en segundo plano del navegador; de este modo, la línea de
comandos y la interfaz web usan el mismo código de análisis y las pruebas
del motor cubren a ambas.

% -------------------------------------------------------------
\section{Evidencia de implementación}\label{sec:evidencia}
% -------------------------------------------------------------

El sistema se encuentra publicado en \url{https://graph-sast.vercel.app/}
y su código fuente en el repositorio indicado en el
anexo~\ref{anx:repositorios}. La versión que corresponde a esta entrega
se identifica con la etiqueta \completar{entrega-7} del repositorio. Las
capturas del sistema en funcionamiento se presentan en el
capítulo~\ref{cap:marco-tecnologico}: la detección de un recorrido
vulnerable, el descarte del mismo recorrido saneado, el panel de
métricas y el análisis de un proyecto cargado como carpeta, en el que el
hallazgo cruza tres archivos.

\completar{Agregar una captura del caso A de la validación con usuarios
cargado en la versión publicada, con el camino entre los dos archivos
resaltado.}

% -------------------------------------------------------------
\section{Pruebas técnicas}\label{sec:pruebas}
% -------------------------------------------------------------

\subsection{Pruebas automatizadas}

La suite automatizada comprende 247 pruebas en 27 archivos, ejecutadas
con Vitest: 219 corresponden al motor y 28 a la interfaz web. La
tabla~\ref{tab:suite} las agrupa por etapa del procesamiento. En la
ejecución del 24 de septiembre de 2026 todas las pruebas resultaron
satisfactorias. La paridad entre el motor Cypher y el recorrido en
memoria (RF-10) no forma parte de la suite, porque requiere una instancia
activa de Neo4j. \completar{Resultado de la verificación de paridad con
Neo4j sobre el banco sintético.}

\begin{table}[H]
    \centering
    \caption{Pruebas automatizadas por etapa}
    \label{tab:suite}
    \small
    \begin{tabularx}{\textwidth}{L r l}
        \toprule
        \tb{Etapa} & \tb{Pruebas} & \tb{Requerimiento} \\
        \midrule
        Análisis sintáctico & 2 & RF-01 \\
        Representación intermedia y grafo de llamadas & 14 & RF-02 \\
        Flujo de datos intra e inter-procedural y ámbitos & 34 & RF-02 \\
        Motor de contaminación y destinos XSS & 26 & RF-03, RF-04 \\
        Catálogo de reglas & 8 & RNF-06 \\
        Veredicto y selección de motor & 18 & RF-03 \\
        Análisis entre archivos & 22 & RF-05 \\
        Escaneo de archivos y exclusiones & 31 & RF-07, RF-08 \\
        Informes de texto, JSON y SARIF & 21 & RF-08 \\
        Códigos de salida & 9 & RF-09 \\
        API local & 19 & RF-11 \\
        Generación de consultas Cypher & 2 & RF-10 \\
        Evaluación y banco de pruebas & 7 & RNF-01 \\
        Punto de entrada del motor & 6 & RF-01 \\
        Interfaz web & 28 & RF-06, RF-07 \\
        \midrule
        \tb{Total} & \tb{247} & \\
        \bottomrule
    \end{tabularx}
\end{table}

\subsection{Banco de pruebas sintético}

\paragraph{Diseño.} El banco se compone de 43 casos de verdad conocida,
29 vulnerables y 14 seguros, con un total de 198 líneas de código. Cada
caso declara si es vulnerable, la familia de debilidad y la cantidad
esperada de hallazgos. Los casos cubren recorridos dentro de una función
y entre funciones, entradas propias de Express (parámetros, cuerpo,
encabezados y cookies), asignaciones, desestructuración, literales de
plantilla, funciones flecha, métodos de clase y componentes JSX. Entre
los casos seguros se incluyen tres construidos contra la detección por
nombres: funciones llamadas \texttt{evaluatePrice}, \texttt{myExecutor}
y \texttt{spawnConfetti}, cuyos nombres contienen los de destinos
sensibles del catálogo (\texttt{eval}, \texttt{exec} y \texttt{spawn})
sin serlo. En sentido inverso, dos casos vulnerables verifican que el
saneamiento se reconoce por su posición en el camino y no por el nombre:
uno aplica el saneamiento en una rama paralela que no protege al destino,
y otro usa una función cuyo nombre contiene \texttt{escape} sin sanear
el dato.

\begin{table}[H]
    \centering
    \caption{Composición del banco sintético por familia}
    \label{tab:banco-composicion}
    \small
    \begin{tabular}{l r r}
        \toprule
        \tb{Familia} & \tb{Vulnerables} & \tb{Seguros} \\
        \midrule
        CWE-89, inyección SQL & 21 & 5 \\
        CWE-79, Cross-Site Scripting & 3 & 3 \\
        CWE-78, inyección de comandos & 1 & 0 \\
        CWE-943, inyección NoSQL & 4 & 0 \\
        Sin familia (nombres engañosos y otros negativos) & 0 & 6 \\
        \midrule
        \tb{Total} & \tb{29} & \tb{14} \\
        \bottomrule
    \end{tabular}
\end{table}

\paragraph{Resultados.} La tabla~\ref{tab:confusion} presenta la matriz
de confusión. El sistema clasificó correctamente los 43 casos: detectó
los 29 vulnerables y no emitió hallazgos sobre los 14 seguros, con una
precisión, una exhaustividad y un $F_1$ de 1,00. El análisis completo
del banco demandó 52\,ms (mediana de cinco ejecuciones), es decir,
0,27\,ms por línea, lo que equivale a 0,27\,s por cada 1.000 líneas y
cumple el umbral de RNF-03.

\begin{table}[H]
    \centering
    \caption{Matriz de confusión sobre el banco sintético (43 casos)}
    \label{tab:confusion}
    \small
    \begin{tabular}{l r r}
        \toprule
        & \tb{Reportado vulnerable} & \tb{No reportado} \\
        \midrule
        \tb{Caso vulnerable} & 29 (TP) & 0 (FN) \\
        \tb{Caso seguro} & 0 (FP) & 14 (TN) \\
        \bottomrule
    \end{tabular}
\end{table}

\paragraph{Alcance del resultado.} El banco fue construido por el autor
y se limita a construcciones que el analizador cubre, de modo que un
resultado perfecto sobre él no constituye evidencia de desempeño sobre
código real. Su función es de regresión: verificar que cada cambio del
motor mantiene la detección de los casos conocidos (RNF-01). La
distribución es además desigual: 26 de los 43 casos corresponden a
inyección SQL y solo uno a inyección de comandos.

\subsection{Comparación con una línea de base}\label{sec:linea-base}

Las hipótesis H1 y H2 comparan el análisis sobre el grafo con la
detección por coincidencia de patrones. Para medir esa diferencia se
aplica sobre los mismos casos una línea de base que
\completar{marca como vulnerable toda llamada cuyo nombre coincide con
un destino del catálogo y que recibe al menos un argumento no literal,
sin seguir el flujo del dato ni reconocer saneamientos}.

\begin{table}[H]
    \centering
    \caption{Resultados de la línea de base y de GraphSAST sobre los mismos casos}
    \label{tab:metricas}
    \small
    \begin{tabular}{l l r r r r r}
        \toprule
        \tb{Corpus} & \tb{Detector} & \tb{FP} & \tb{FN} & \tb{Precisión} & \tb{Exhaustividad} & \tb{$F_1$} \\
        \midrule
        Banco sintético & Línea de base & \completar{..} & \completar{..} & \completar{..} & \completar{..} & \completar{..} \\
        Banco sintético & GraphSAST & 0 & 0 & 1,00 & 1,00 & 1,00 \\
        NodeGoat & Línea de base & \completar{..} & \completar{..} & \completar{..} & \completar{..} & \completar{..} \\
        NodeGoat & GraphSAST & \completar{..} & \completar{..} & \completar{..} & \completar{..} & \completar{..} \\
        \bottomrule
    \end{tabular}
\end{table}

\completar{Lectura de la comparación: en qué casos difieren los dos
detectores y por qué.}

\subsection{Vulnerabilidades documentadas en paquetes de npm}

Para contrastar el sistema con vulnerabilidades reales se analizaron seis
advisories del ecosistema npm publicados en la base de datos de GitHub,
relevados el 1 de septiembre de 2026 y congelados para que la medición
sea reproducible. Para cada caso se analizaron la versión vulnerable y la
corregida del paquete, contando solo los hallazgos de la familia del
advisory. El detalle se presenta en el
capítulo~\ref{cap:marco-tecnologico}.

El sistema detectó uno de los seis casos (16,7\,\%) y ninguno con
discriminación: el único hallazgo, en DOMPurify, se mantiene en la
versión corregida, por lo que no corresponde a la vulnerabilidad
reportada. En dos paquetes el registro de npm solo contiene código
compilado; en otro, el paquete publicado es un instalador de tres
archivos; y en los restantes la falla reside en la lógica interna de
filtrado de bibliotecas de saneamiento, que no es un recorrido de un
origen a un destino. El corpus, por lo tanto, no permitió evaluar el
modelo de amenaza del analizador, y este resultado dio lugar a la
respuesta al riesgo R-02: sumar una aplicación deliberadamente
vulnerable.

\subsection{Aplicación deliberadamente vulnerable: OWASP NodeGoat}

\paragraph{Diseño.} OWASP NodeGoat es una aplicación Node.js construida
con fines didácticos para mostrar los riesgos del OWASP Top 10,
distribuida con licencia Apache-2.0. Antes de ejecutar el analizador se
etiquetan las vulnerabilidades documentadas que corresponden a las
familias del catálogo, de modo que la medición no se ajuste a los
resultados. \completar{Versión o commit analizado y fecha.}

\paragraph{Resultados.} \completar{Cantidad de vulnerabilidades
etiquetadas por familia, hallazgos, matriz de confusión y exhaustividad
desglosada por construcción sintáctica (riesgo R-01).}

\subsection{Rendimiento}

Sobre el banco sintético el análisis demandó 0,27\,ms por línea en la
línea de comandos. En la versión web, el análisis de un proyecto de 78
archivos y unas 8.500 líneas demandó alrededor de 2,8\,s en el equipo de
desarrollo.
\completar{Tiempo de análisis de NodeGoat en la línea de comandos y en
la versión web.}

% -------------------------------------------------------------
\section{Validación con usuarios y con un experto}\label{sec:validacion-usuarios}
% -------------------------------------------------------------

\paragraph{Diseño.} La hipótesis H3 se evalúa con un estudio
intra-sujeto: cada participante resuelve un caso con el reporte de texto
de la línea de comandos y otro con la visualización del grafo, y el
orden de las condiciones y de los casos se alterna entre participantes.
Los dos casos tienen la misma estructura, dos archivos con un camino
vulnerable que los cruza y un endpoint saneado, y difieren en la
debilidad: inyección SQL en el caso A e inyección de comandos en el caso
B. Por cada tarea se registran el acierto en cuatro preguntas, el tiempo
y dos valoraciones en escala de 1 a 5. El protocolo completo y el
cuestionario se presentan en el apéndice~\ref{apx:encuesta}. De forma
complementaria se entrevistó a un desarrollador con experiencia en
seguridad de aplicaciones (apéndice~\ref{apx:entrevista}).

\paragraph{Participantes.} \completar{Cantidad, experiencia con
JavaScript, período de las sesiones.}

\begin{table}[H]
    \centering
    \caption{Resultados por condición en la validación con usuarios}
    \label{tab:h3-resultados}
    \small
    \begin{tabular}{l r r}
        \toprule
        \tb{Indicador} & \tb{Reporte de texto} & \tb{Grafo} \\
        \midrule
        Acierto, mediana (0 a 4) & \completar{..} & \completar{..} \\
        Tiempo, mediana (s) & \completar{..} & \completar{..} \\
        L1 comprensión del recorrido, mediana (1 a 5) & \completar{..} & \completar{..} \\
        L2 claridad de la corrección, mediana (1 a 5) & \completar{..} & \completar{..} \\
        Participantes con mejor resultado en la condición & \completar{..} & \completar{..} \\
        \bottomrule
    \end{tabular}
\end{table}

\completar{Preferencia declarada, respuestas sobre la disposición a pagar
USD 8 mensuales (riesgo R-07) y observaciones de las sesiones.}

\paragraph{Entrevista.} \completar{Síntesis de la entrevista por tema.}

\begin{table}[H]
    \centering
    \caption{Traducción de la retroalimentación en decisiones de diseño}
    \label{tab:feedback-decisiones}
    \small
    \begin{tabularx}{\textwidth}{l L L l}
        \toprule
        \tb{Fuente} & \tb{Hallazgo} & \tb{Decisión} & \tb{Estado} \\
        \midrule
        Sesiones (n=\completar{..}) & \completar{...} & \completar{...} & \completar{...} \\
        Entrevista & \completar{...} & \completar{...} & \completar{...} \\
        \bottomrule
    \end{tabularx}
\end{table}

% -------------------------------------------------------------
\section{Resultados desfavorables y su análisis}\label{sec:resultados-negativos}
% -------------------------------------------------------------

\paragraph{Validación sobre paquetes de npm.} Como se expuso en la
sección~\ref{sec:pruebas}, solo uno de seis advisories produjo un
hallazgo, y ese hallazgo no discrimina entre la versión vulnerable y la
corregida. El resultado no confirma ni refuta la efectividad del sistema,
porque las fallas del corpus no son recorridos de un origen a un
destino, pero deja sin evidencia sobre código real a la validación
prevista en el capítulo~\ref{cap:metodologia} hasta contar con los
resultados de NodeGoat.

\paragraph{Falsos positivos por coincidencia de nombres.} Al analizar el
código del propio sistema, una llamada a \texttt{RegExp.prototype.exec}
se reportó como inyección de comandos, porque coincide por nombre con el
destino \texttt{exec} del catálogo. El banco sintético no detectó el
problema porque sus casos engañosos usan nombres que contienen al de un
destino, como \texttt{myExecutor}, y no el mismo nombre invocado sobre
otro objeto. \completar{Estado de la corrección: calificar el destino por
el módulo del que se importa (riesgo R-04).}

\paragraph{Cobertura del análisis entre archivos.} El dato se sigue entre
archivos solo cuando la función invocada se declara con
\texttt{function} y se importa con módulos ES y ruta relativa. Dentro de
un mismo archivo, en cambio, el banco sintético incluye casos con
funciones flecha y con un método de clase que se detectan correctamente.
La limitación afecta, por lo tanto, a las aplicaciones Express
organizadas en varios archivos con funciones flecha o con
\texttt{require}, que es la forma habitual de NodeGoat. \completar{Efecto
medido sobre NodeGoat.}

% -------------------------------------------------------------
\section{Síntesis de la evidencia}\label{sec:sintesis-evidencia}
% -------------------------------------------------------------

La tabla~\ref{tab:evidencia-hipotesis} resume qué evidencia se obtuvo
para cada hipótesis. La respuesta a cada una se desarrolla en el
capítulo~\ref{cap:conclusiones}.

\begin{table}[H]
    \centering
    \caption{Evidencia obtenida por hipótesis}
    \label{tab:evidencia-hipotesis}
    \small
    \begin{tabularx}{\textwidth}{l L L}
        \toprule
        \tb{Hipótesis} & \tb{Evidencia obtenida} & \tb{Sección} \\
        \midrule
        H1 & Banco sintético: 29 de 29 casos vulnerables detectados, incluidos recorridos entre funciones y entre archivos. \completar{Comparación con la línea de base y resultados de NodeGoat.} & \ref{sec:pruebas} \\
        H2 & Banco sintético: 0 falsos positivos en 14 casos seguros, tres de ellos saneados y tres con nombres engañosos; dos casos vulnerables con saneamientos que no protegen el destino se detectan. Un falso positivo por coincidencia de nombres fuera del banco. \completar{Falsos positivos de la línea de base.} & \ref{sec:pruebas}, \ref{sec:resultados-negativos} \\
        H3 & \completar{Resultados de las sesiones y de la entrevista.} & \ref{sec:validacion-usuarios} \\
        \bottomrule
    \end{tabularx}
\end{table}
